%% guide-a1.tex　ガイドブック　保護者・先生のための本（大人用、vol=0）第1章「シリーズの概要と使い方」
%% （第16回編集委員会の追加の決定、組 G-B、頭文字 kgd）。ルビなし。大人向けの漢字の使い方（10月30日 10:15 依頼者）。
%% 材料：tools/guide-series.md（RG-2）。数・事実は材料のとおり。
%% このファイルの先頭で、大人用の本文（guide-a1・guide-p2a・guide-p2a-src）と付録（guide-p2b、G-C）が共通に使う部品を定義する。
%% 親ファイル（tex/guide-adult.tex、G-A）は guide-a1 → guide-p2a → guide-p2b → guide-p2a-src の順に読む。
\protected\def\R#1#2{#1}% 大人用の本はルビなし（親ファイルでも同じ定義をする）
% 文中の「→」のあとで改行しない（行末に→を残さない。大人用の本の全体に効く）
\ltjsetparameter{postbreakpenalty={`→,10000}}
% ★の並び（★★〜★★★ など）が行末で切れないようにする（委員確認 rA #22。付録にも効く）
\ltjsetparameter{postbreakpenalty={`★,10000}}\ltjsetparameter{postbreakpenalty={`☆,10000}}
\ltjsetparameter{prebreakpenalty={`〜,10000}}\ltjsetparameter{postbreakpenalty={`〜,10000}}
\makeatletter
%% ===== 章見出し・小見出し（guide-p2a・guide-p2b・guide-p2a-src と共通。G-C は \providecommand の写しを持つ） =====
% \kgbsec[柱・しおり用の題（省略可）]{番号}{題}。改ページはしない（前の章に続けて組む。見出しだけがページの下に残らないよう 8 行分あける）
\providecommand\kgbsec[3][]{%
  \par\addvspace{18pt}\Needspace{8\baselineskip}\global\bk@decotrue
  \setcounter{bkfig}{0}\setcounter{bktab}{0}\gdef\kgbsecno{#2}%
  \if\relax\detokenize{#1}\relax\def\kgb@st{#3}\else\def\kgb@st{#1}\fi
  \bksetrunning{}{保護者・先生のための本}{\if\relax\detokenize{#2}\relax\else#2\hspace{0.6\zw}\fi\kgb@st}%
  \pdfbookmark[1]{#2　\kgb@st}{kgbsec-#2}%
  \kgb@secband{#2}{#3}}
\newcommand\kgb@secband[2]{%
  \noindent\begin{tikzpicture}
    \if\relax\detokenize{#1}\relax
      \node[anchor=west, text=white, font=\sffamily\bfseries\LARGE, inner sep=0pt] (t) at (0.6,0) {#2};
    \else
      \node[anchor=west, circle, fill=white, text=bkVol0, minimum size=1.35cm, font=\sffamily\bfseries\Large, inner sep=0pt] (n) at (0.35,0) {#1};
      \node[anchor=west, text=white, font=\sffamily\bfseries\LARGE, text width=\linewidth-2.4cm, align=left, inner sep=0pt] (t) at ([xshift=0.4cm]n.east) {#2};
    \fi
    \begin{pgfonlayer}{background}
      \path let \p1=(t.north), \p2=(t.south) in coordinate (kgbT) at (0,{max(\y1+0.4cm,1cm)}) coordinate (kgbB) at (0,{min(\y2-0.4cm,-1cm)});
      \fill[bkVol0] (kgbB) rectangle (\linewidth,0 |- kgbT);
    \end{pgfonlayer}
  \end{tikzpicture}\par\vspace{10pt}}
\providecommand\kgbsub[2]{%
  \par\addvspace{12pt}\Needspace{5\baselineskip}\noindent
  \begin{tikzpicture}[baseline=(t.base)]
    \node[anchor=west, fill=bkVol0, text=white, rounded corners=1mm, font=\sffamily\bfseries\small, inner xsep=1.6mm, inner ysep=1mm] (i) at (0,0) {#1};
    \node[anchor=west, font=\sffamily\bfseries\large, text=bkVol0, inner sep=0pt] (t) at ([xshift=0.6\zw]i.east) {#2};
    \draw[bkVol0, line width=0.8pt] ([yshift=-4pt]i.south west) -- ++(\linewidth-1pt,0);
  \end{tikzpicture}\par\nobreak\vspace{6pt}}
% 図・表の番号は章ごとに「章-番号」（図4-1、表6-2 など）
\providecommand\kgbsecno{}
\renewcommand\thebkfig{\kgbsecno-\arabic{bkfig}}\renewcommand\thebktab{\kgbsecno-\arabic{bktab}}
\renewcommand\theHbkfig{kgb.\kgbsecno.\arabic{bkfig}}\renewcommand\theHbktab{kgb.\kgbsecno.\arabic{bktab}}% リンクの名前がぶつからないように
%% ===== 大人用の本文（第1〜10章）で使う部品（頭文字 kgb。名前は第2部のときのまま） =====
% 巻の札  \kgbvol{1〜9}（巻の色の地に、その巻の帯の字の色）
\newcommand\kgbvol[1]{\tikz[baseline=(v.base)]\node[fill=bkVol#1, text=bkOn#1, rounded corners=0.8mm, inner xsep=1.2mm, inner ysep=0.7mm, font=\sffamily\bfseries\footnotesize] (v) {第#1巻};}
% 箇条書き（1字下げのぶら下がり）  \kgbitem{文}
\newcommand\kgbitem[1]{\par\noindent\hangindent=1\zw\hangafter=1 ・#1\par}
% 注（表・図の下の小さな字の説明）  \kgbnote{文}
\newcommand\kgbnote[1]{\par\vspace{1pt}{\footnotesize\noindent\hangindent=1\zw\hangafter=1 #1\par}\vspace{2pt}}
% ★の範囲（表の中で）  \kgbst{最小}{最大}
\newcommand\kgbst[2]{\mbox{\footnotesize\bkstars{#1}\ifnum#1<#2\relax{\color{bkInk}〜}\bkstars{#2}\fi}}
% 長い表（ページをまたげる表。続きのページでも1行めの見出しをくり返す）
% 表の名前は、表の最後の行（線なし）に置き、すぐ上の行とページで離れないようにする（\\*）。
% 使い方：\kgbcap{表の名前}{ラベル} を \begin{bktable}[long, theme=kgb] の直前に置く
\NewTblrTheme{kgb}{\DefTblrTemplate{caption}{default}{}\DefTblrTemplate{conthead}{default}{}%
  \DefTblrTemplate{capcont}{default}{}\DefTblrTemplate{contfoot}{default}{}%
  \DefTblrTemplate{note}{default}{}\DefTblrTemplate{remark}{default}{}}
\def\kgb@capt{}
\newcommand\kgbcap[2]{\refstepcounter{bktab}\if\relax\detokenize{#2}\relax\else\label{#2}\fi\gdef\kgb@capt{#1}}
\newcommand\kgbcapline{{\sffamily\small{\color{\bkTT}\bfseries 表\thebktab}\hspace{0.8\zw}\kgb@capt}}
% 学年と巻の対応の帯の図：左の巻の欄 24 mm、小1〜小6・中1〜中3・高校の6科目は 7.2 mm、いちばん右の「大学」の列は 8 mm
\newcommand\kgbgridcell[3]{% {y}{巻}{列} 内容がある学年・科目のます
  \pgfmathsetmacro\kgb@x{ifthenelse(#3<16,24+(#3-1)*7.2,132)}\pgfmathsetmacro\kgb@w{ifthenelse(#3<16,7.2,8)}%
  \fill[bkVol#2, draw=bkInk!55, line width=0.4pt] ({\kgb@x+(\kgb@w-5.4)/2},{#1-2.7}) rectangle ++(5.4,5.4);}
\newcommand\kgbgridband[4]{% {y}{巻}{はじめの列}{終わりの列} 学年の目安の太い枠
  \pgfmathsetmacro\kgb@xa{ifthenelse(#3<16,24+(#3-1)*7.2,132)}\pgfmathsetmacro\kgb@xb{ifthenelse(#4<16,24+#4*7.2,140)}%
  \draw[bkInk, line width=1.3pt, rounded corners=1.3mm] ({\kgb@xa+0.35},{#1-3.6}) rectangle ({\kgb@xb-0.35},{#1+3.6});}
\newcommand\kgbgridempty[1]{% {y} 空のますのうすい枠（16列）
  \foreach \kgb@c in {1,...,16} {\pgfmathsetmacro\kgb@x{ifthenelse(\kgb@c<16,24+(\kgb@c-1)*7.2,132)}\pgfmathsetmacro\kgb@w{ifthenelse(\kgb@c<16,7.2,8)}%
    \draw[bkLine, line width=0.4pt] ({\kgb@x+(\kgb@w-5.4)/2},{#1-2.7}) rectangle ++(5.4,5.4);}}
% 1行  \kgbgridrow{y}{巻}{内容がある列のならび}{目安のはじめの列}{目安の終わりの列}
\newcommand\kgbgridrow[5]{%
  \fill[bkPaper] (132,{#1-4.1}) rectangle (140,{#1+4.1});
  \kgbgridempty{#1}%
  \foreach \kgb@c in {#3} {\kgbgridcell{#1}{#2}{\kgb@c}}%
  \kgbgridband{#1}{#2}{#4}{#5}%
  \node[anchor=west, inner sep=0pt] at (0.5,#1) {\kgbvol{#2}};}
% 見出し（上から y=0 まで）。#1＝高校の上に学年の例の行をつけるか（1/0）
\newcommand\kgbgridhead[1]{%
  \tikzset{kgb gh/.style={font=\sffamily\bfseries\footnotesize, text=bkInk, inner sep=0.5pt}}
  \ifnum#1=1 \def\kgb@hy{12}\else\def\kgb@hy{4.8}\fi
  \fill[bkPaper] (132,0) rectangle (140,\kgb@hy+7);
  \node[kgb gh] at ({24+3*7.2},\kgb@hy+3.5) {小学校（学年）};
  \node[kgb gh] at ({24+7.5*7.2},\kgb@hy+3.5) {中学校};
  \node[kgb gh] at ({24+12*7.2},\kgb@hy+3.5) {高校（数学の科目）};
  \node[kgb gh] at (136,\kgb@hy+3.5) {大学};
  \foreach \kgb@c/\kgb@t in {1/1,2/2,3/3,4/4,5/5,6/6,7/1,8/2,9/3,10/Ⅰ,11/A,12/Ⅱ,13/B,14/Ⅲ,15/C}
    {\node[kgb gh] at ({24+(\kgb@c-0.5)*7.2},3) {\kgb@t};}
  \ifnum#1=1
    \foreach \kgb@c/\kgb@t in {10/1年の例,12/2年の例,14/3年の例}
      {\node[kgb gh, font=\sffamily\footnotesize] at ({24+\kgb@c*7.2},8.6) {\kgb@t};
       \draw[bkGray, line width=0.5pt] ({24+(\kgb@c-1)*7.2+0.8},6.6) -- ({24+(\kgb@c+1)*7.2-0.8},6.6);}
  \fi
  \foreach \kgb@x in {24,{24+6*7.2},{24+9*7.2},132,140} {\draw[bkGray, line width=0.6pt] (\kgb@x,\kgb@hy+7) -- (\kgb@x,0);}
  \draw[bkGray, line width=0.6pt] (24,0) -- (140,0);}
% 縦の区切り線を下まで  \kgbgridlines{いちばん下の y}
\newcommand\kgbgridlines[1]{\foreach \kgb@x in {24,{24+6*7.2},{24+9*7.2},132,140} {\draw[bkGray, line width=0.6pt] (\kgb@x,0) -- (\kgb@x,#1);}}
% 各巻のカード（5章）の1行の帯  \kgbstrip{巻}{内容がある列}{目安のはじめ}{目安の終わり}
\newcommand\kgbstrip[4]{\begin{tikzpicture}[x=0.935mm, y=1mm]
  \kgbgridhead{0}\kgbgridrow{-5}{#1}{#2}{#3}{#4}\kgbgridlines{-9.1}\end{tikzpicture}}
% 巻のカード（5章）  \begin{kgbcard}{巻}{題}
\newtcolorbox{kgbcard}[2]{enhanced, breakable=false, arc=2mm, boxrule=1pt, colframe=bkVol#1, colback=white,
  left=3mm, right=3mm, top=1.2mm, bottom=1.6mm, before skip=6pt, after skip=6pt,
  title={\kgbvol{#1}\hspace{0.6\zw}{\color{bkInk}#2}}, colbacktitle=bkVol#1!14, coltitle=bkInk,
  fonttitle=\sffamily\bfseries, toptitle=0.6mm, bottomtitle=0.6mm, titlerule=0pt}
% カードの中の1行（見出し＋文）  \kgbrow{見出し}{文}
\newcommand\kgbrow[2]{\par\addvspace{1pt}{\footnotesize\baselineskip=15.2pt\noindent\hangindent=6.6\zw\hangafter=1
  \makebox[6.6\zw][l]{\sffamily\bfseries\color{bkVol0}#1}#2\par}}
% 図の箱
\tikzset{
  kgb box/.style={draw=bkVol0, line width=0.9pt, fill=white, rounded corners=1.5mm, align=center,
    font=\sffamily\footnotesize, text=bkInk, inner xsep=2mm, inner ysep=1.6mm},
  kgb gbox/.style={draw=bkGray, line width=0.8pt, fill=white, rounded corners=1.5mm, align=left,
    font=\footnotesize, text=bkInk, inner xsep=1.8mm, inner ysep=1.4mm},
  kgb tbox/.style={draw=bkVol0, line width=0.9pt, fill=bkVol0!8, rounded corners=1.5mm, align=left,
    font=\footnotesize, text=bkInk, inner xsep=1.8mm, inner ysep=1.4mm},
  kgb arr/.style={-{Stealth[length=2.4mm]}, line width=1pt, bkInk!75, shorten <=1.2mm, shorten >=1.2mm},
  kgb lab/.style={font=\sffamily\footnotesize, text=bkInk, inner sep=0.8pt},
}
\makeatother

%% ===== 第1章だけで使う部品（頭文字 kgd） =====
\makeatletter
% 本棚の1冊（背）  \kgdspine{x}{巻}{はば}{題}{ページ}
\newcommand\kgdspine[5]{%
  \fill[bkVol#2, draw=bkInk!60, line width=0.5pt, rounded corners=0.6mm] (#1,0) rectangle ++(#3,64);
  \node[text=bkOn#2, font=\sffamily\bfseries\footnotesize, inner sep=0pt] at ({#1+#3/2},58.5) {#2};
  \node[text=bkOn#2, font=\sffamily\bfseries\footnotesize, inner sep=0pt, anchor=north, align=center, text width=#3 mm, execute at begin node={\baselineskip=10.6pt}] at ({#1+#3/2},53.5) {#4};% 題は1字ずつ縦に積む（\\ で区切る）
  \node[font=\sffamily\footnotesize, text=bkInk, inner sep=0pt, anchor=north] at ({#1+#3/2},-1.6) {#5};}
% 章の流れの1つの札  \kgdpart{名前}{アイコン}{x}{y}{ノード名}
\newcommand\kgdpart[5]{\node[kgb box, anchor=west, minimum width=43.5mm, text width=39.5mm, minimum height=8.6mm, align=left, inner xsep=1.5mm, inner ysep=1mm] (#5) at (#3,#4)
  {\makebox[5mm][l]{\if\relax\detokenize{#2}\relax\else\ifcsname kgd@ibox@#2\endcsname\usebox{\csname kgd@ibox@#2\endcsname}\else\bkicon[1.1em]{#2}{bkVol0}\fi\fi}\hspace{0.5mm}#1};}
% 中に \node を持つアイコン（はてな）は、図の外で箱に入れておく（図の中で入れ子にすると位置がずれるため。委員確認 rA #20）
\newsavebox\kgd@ibox@hatena
\sbox\kgd@ibox@hatena{\bkicon[1.1em]{hatena}{bkVol0}}
\makeatother

%% ====================================================================
\kgbsec{1}{シリーズの概要と使い方}\label{kgd-ch1}

この本は、『物理を理解するための数学』全9巻を使う保護者・先生のための案内です。第1章でシリーズの概要と使い方を、第2章以降で編集方針と学習指導要領との対応を説明します。子ども向けの『はじめに読む本』を読まなくても、この1冊でシリーズの全体と使い方が分かるように書きました。

\kgbsub{1.1}{このシリーズの目的}

このシリーズは、工学の研究を発展させるために、小学校から高校までの算数・数学の学び方を見直す、という考えのもとで作られました。小学1年の「リンゴは何個？」から始めて、工学部で使う行列・微分方程式・ベクトル解析・フーリエ解析・ラプラス変換までを、一続きの9巻で学べるようにしています。

狙いは、入試を解くための発想力の数学ではなく、物理や工学の現象を説明するための数学を、子どもの「なぜ？」「何がうれしいの？」に答えながら教えることです。このシリーズを終えてから大学の専門書を開くと、書いてあることがすんなり頭に入る状態を目標にしています。

\begin{bkfigure}
\begin{tikzpicture}[x=1mm, y=1mm]
  \node[kgb box, text width=30mm] (a) at (17,0) {第1巻 1-1\\[1pt]{\bfseries いくつ あるかな？}\\（数える・比べる）};
  \node[kgb box, text width=36mm, minimum height=19mm] (b) at (70,0) {7本の路線\\[1pt]{\mdseries 変化・たし合わせ・倍々・\\回転と波・表・矢印・はかる}};
  \node[kgb box, text width=32mm, fill=bkVol0!8] (c) at (120.5,0) {工学部の専門書\\[1pt]{\mdseries 行列・微分方程式・\\ベクトル解析・\\フーリエ解析・\\ラプラス変換}};
  \draw[kgb arr] (a) -- (b);
  \draw[kgb arr] (b) -- (c);
  \node[kgb lab, anchor=north] at (70,-21) {全9巻・152章（4836ページ）を、7本の路線でつなぐ};
\end{tikzpicture}
\bkfigcaption{このシリーズの出発点と到達点}
\end{bkfigure}

本文の表記は、読む子どもの学年に合わせています。第1〜3巻は全ての漢字にルビを付け（第1巻は分かち書きも付けます）、第4巻からはルビを付けずに常用漢字までを使います。

\kgbsub{1.2}{9巻の並び}

全9巻は、数えることから始めて、関数とデータ、図形とベクトル、微分と積分を経て、大学の工学の数学へ進みます。

\begin{bkfigure}
\begin{tikzpicture}[x=1mm, y=1mm]
  \draw[bkGray, line width=1.2pt] (0.6,-0.3) -- (139.4,-0.3);
  \kgdspine{0.10}{1}{11.11}{か\\ぞ\\え\\る\\・\\は\\か\\る}{400}
  \kgdspine{11.89}{2}{12.67}{数\\と\\量\\の\\こ\\と\\ば}{456}
  \kgdspine{25.24}{3}{15.67}{関\\数\\と\\デ\\\rotatebox[origin=c]{-90}{ー}\\タ}{564}
  \kgdspine{41.59}{4}{15.56}{形\\と\\向\\き}{560}
  \kgdspine{57.83}{5}{13.67}{変\\化\\を\\つ\\か\\む}{492}
  \kgdspine{72.18}{6}{16.67}{未\\来\\を\\計\\算\\す\\る}{600}
  \kgdspine{89.53}{7}{13.67}{ま\\と\\め\\て\\計\\算\\す\\る}{492}
  \kgdspine{103.88}{8}{17.56}{空\\間\\に\\広\\が\\る\\量}{632}
  \kgdspine{122.12}{9}{17.78}{波\\と\\変\\化\\を\\読\\み\\と\\く}{640}
  \node[kgb lab, anchor=north west, text=bkGray] at (0,-6) {本の下の数はページ数。背の幅はページ数に比例させた（1 mm＝36ページ）};
\end{tikzpicture}
\bkfigcaption{全9巻の本棚（色は各巻の表紙の色）}
\end{bkfigure}

\par\addvspace{4pt}
\kgbcap{9巻の一覧（学年の目安・★の範囲・章の数・ページ数）}{kgd-t-vols}
\begin{bktable}[long, theme=kgb]{hline{Z}={0pt,white}, vline{1-Z}={Z}{0pt,white}, row{Z}={abovesep=3pt, belowsep=1pt}, hspan=minimal, rowhead=1, colspec={Q[c,m,wd=13mm] Q[l,m,wd=35mm] Q[l,m,wd=26mm] Q[c,m,wd=31mm] Q[c,m,wd=8mm] Q[r,m,wd=11mm]}, rowsep=1.5pt, colsep=3pt, cells={font=\small}, row{1}={font=\sffamily\bfseries\small}}
巻 & 書名（副題） & 学年の目安 & ★の範囲 & 章 & ページ \\
\kgbvol{1} & かぞえる・はかる & 小学1年〜3年 & \kgbst{1}{1} & 22 & 400 \\
\kgbvol{2} & 数と量のことば & 小学4年〜中学1年 & \kgbst{1}{2} & 20 & 456 \\
\kgbvol{3} & 関数とデータ（きまりを見つける） & 中学1年〜3年 & \kgbst{1}{3} & 18 & 564 \\
\kgbvol{4} & 形と向き（図形・三角関数・ベクトル） & 中学2年〜高校2年 & \kgbst{1}{5} & 16 & 560 \\
\kgbvol{5} & 変化をつかむ（微分と積分） & 高校1年〜2年 & \kgbst{2}{4} & 19 & 492 \\
\kgbvol{6} & 未来を計算する（微分方程式と複素数） & 高校2年〜3年 & \kgbst{2}{5} & 16 & 600 \\
\kgbvol{7} & まとめて計算する（行列） & 高校3年（発展）〜大学1年 & \kgbst{3}{5} & 12 & 492 \\
\kgbvol{8} & 空間に広がる量（ベクトル解析） & 大学1年〜2年 & \kgbst{3}{5} & 15 & 632 \\
\kgbvol{9} & 波と変化を読みとく（フーリエ解析とラプラス変換） & 大学2年 & \kgbst{3}{5} & 14 & 640 \\
計 & & & & 152 & 4836 \\*
\SetCell[c=6]{c} \kgbcapline & & & & & \\
\end{bktable}
\par\addvspace{4pt}

\kgbnote{※ 章の数は序章を含みます。第7・8巻は読みもの1章、第9巻は読みもの1章と終章を含みます（全152章＝序章9・章139・読みもの3・終章1）。ページ数は奥付までのページです。}
\kgbnote{※ 学年の目安は、編集のときに決めた目安で、生徒用の紙面には書いていません（理由と読み方は 4.2・9.3）。実際には、目安の学年より前や後に学校で習う内容も含まれます（図\ref{kgb-f-grid}）。}

各巻は、いくつかの「部」に分かれています。

\par\addvspace{4pt}
\kgbcap{各巻の部の構成}{}
\begin{bktable}[long, theme=kgb]{hline{Z}={0pt,white}, vline{1-Z}={Z}{0pt,white}, row{Z}={abovesep=3pt, belowsep=1pt}, hspan=minimal, rowhead=1, colspec={Q[c,m,wd=13mm] Q[l,m,wd=113mm]}, rowsep=1.2pt, colsep=4pt, cells={font=\small}, row{1}={font=\sffamily\bfseries\small}}
巻 & 部の構成（序章の後の部の名前） \\
\kgbvol{1} & かずの はじまり／はかる／かけ算と わり算／表・グラフ・かたち \\
\kgbvol{2} & 量と数／大きな数・小さな数・単位／文字と式／コンピュータで計算 \\
\kgbvol{3} & 関数とグラフ／曲がったグラフ／倍々とけた／データと統計 \\
\kgbvol{4} & 図形／三角比と三角関数／ベクトル \\
\kgbvol{5} & 微分／積分／特別な関数 \\
\kgbvol{6} & 微分方程式の入口／複素数と双曲線関数／式で解く \\
\kgbvol{7} & 行列の計算／回転と座標の変換／固有値と行列の使いみち \\
\kgbvol{8} & たくさんの変数の微分／ベクトル場と積分／発散と回転／偏微分方程式 \\
\kgbvol{9} & フーリエ解析／ラプラス変換／行列と微分方程式／終章 \\*
\SetCell[c=2]{c} \kgbcapline & \\
\end{bktable}
\par\addvspace{4pt}

\kgbsub{1.3}{7本の路線}

全9巻の内容は、7本の「路線」でつながっています。どの章も、どれかの路線の駅になっています。①〜④の路線は、第6巻の微分方程式で乗り換えます。各巻の前付けには全9巻の路線図があり（図\ref{kgd-f-map}）、第9巻の終章（9-13）では、7本の路線がどこで合流したかをたどります。路線の色には、必ず形の記号と番号を添えています（第1巻の紙面では路線名を「たしあわせ」「やじるし」と平仮名で書きます）。

\par\addvspace{4pt}\Needspace{80mm}% 表1-3 を1ページに
\kgbcap{7本の路線（主な駅は巻-章の番号）}{kgd-t-route}
\begin{bktable}[long, theme=kgb]{hline{Z}={0pt,white}, vline{1-Z}={Z}{0pt,white}, row{Z}={abovesep=3pt, belowsep=1pt}, hspan=minimal, rowhead=1, colspec={Q[l,m,wd=27mm] Q[l,m,wd=31mm] Q[l,m,wd=49mm] Q[l,m,wd=22mm]}, rowsep=1.2pt, colsep=3pt, cells={font=\footnotesize}, row{1}={font=\sffamily\bfseries\footnotesize}}
路線 & 流れ & 主な駅（巻-章） & 終点 \\
\bkrouteicon[1.1em]{1}\ ① 変化 & ひき算→傾き→微分 & 1-5・2-3・3-5・5-1・6-1・8-4 & 流体・熱 \\
\bkrouteicon[1.1em]{2}\ ② たし合わせ & かけ算→面積→積分 & 1-21・2-3・4-5・5-10・6-3・8-11 & 電磁気 \\
\bkrouteicon[1.1em]{3}\ ③ 倍々 & 累乗→指数→$e$ & 2-7・3-11・5-14・6-4・9-5 & 制御 \\
\bkrouteicon[1.1em]{4}\ ④ 回転と波 & 角→sin→複素数 & 1-20・4-8・5-15・6-10・9-1 & 音・交流 \\
\bkrouteicon[1.1em]{5}\ ⑤ 表 & 九九の表→行列 & 1-13・2-18・3-6・7-1・9-10 & 回路網・制御 \\
\bkrouteicon[1.1em]{6}\ ⑥ 矢印 & 向き→ベクトル→場 & 2-2・3-2・4-10・7-7・8-5 & ロボット・電磁気 \\
\bkrouteicon[1.1em]{7}\ ⑦ はかる & 目もり→誤差→統計 & 1-11・2-6・3-14・9-4 & 計測・診断 \\*
\SetCell[c=4]{c} \kgbcapline & & & \\
\end{bktable}
\par\addvspace{4pt}


\kgbsub{1.4}{1つの章の構成と紙面の部品}

第2巻からは、どの章も同じ順番で進みます（第1巻も、言葉を易しくした同じ形です）。章の長さは章ごとに違います（序章を除くと、短い章で14ページ、長い章で54ページ）。章と部扉は、必ず偶数ページ（見開きの左）から始まります。

\begin{bkfigure}
\begin{tikzpicture}[x=1mm, y=1mm]
  \foreach \kgdn/\kgdi/\kgdx/\kgdy/\kgdk in {
    章の帯・道しるべ//0/0/p1, できるようになること/goal/0/-11/p2, はてな？/hatena/0/-22/p3, やってみよう/yattemiyou/0/-33/p4,
    思い出そう/omoidasu/0/-44/p5, 数学の道具/doguu/0/-55/p6,
    もっと知りたい・\\せいかくに言うと/motto/47.5/0/p7, たいせつ/taisetsu/47.5/-11/p8, ことば・記号カード/card/47.5/-22/p9,
    例題・じぶんでとこう/reidai/47.5/-33/p10, まちがえやすいところ/tsumazuki/47.5/-44/p11, 物理・工学では\\こう使う/butsuri/47.5/-55/p12,
    章末問題/jibun/95/0/p13, ふりかえり/check/95/-11/p14} {\kgdpart{\kgdn}{\kgdi}{\kgdx}{\kgdy}{\kgdk}}
  \foreach \kgda/\kgdb in {p1/p2,p2/p3,p3/p4,p4/p5,p5/p6,p7/p8,p8/p9,p9/p10,p10/p11,p11/p12,p13/p14} {\draw[kgb arr, shorten <=0.4mm, shorten >=0.4mm] (\kgda) -- (\kgdb);}
  \draw[kgb arr, rounded corners=2mm] (p6.east) -- ++(2.5,0) |- (p7.west);
  \draw[kgb arr, rounded corners=2mm] (p12.east) -- ++(2.5,0) |- (p13.west);
  \node[kgb gbox, anchor=north west, text width=39.5mm] at (95,-20) {{\sffamily\bfseries 専門書をのぞいてみよう}\\★★★★以上の章。専門書での書き方を1行ずつ訳す};
  \node[kgb gbox, anchor=north west, text width=39.5mm] at (95,-40) {{\sffamily\bfseries ふきだし}\\案内役が次の章へつなぐ（1見開きに1回まで）};
\end{tikzpicture}
\bkfigcaption{1つの章の流れ（左の列から順に。右下の2つは、章の途中に出てくる部品）}\label{kgd-f-flow}
\end{bkfigure}

\par\addvspace{4pt}
\kgbcap{紙面の部品と役割}{kgd-t-parts}
\begin{bktable}[long, theme=kgb]{hline{Z}={0pt,white}, vline{1-Z}={Z}{0pt,white}, row{Z}={abovesep=3pt, belowsep=1pt}, hspan=minimal, rowhead=1, colspec={Q[l,m,wd=36mm] Q[l,m,wd=92mm]}, rowsep=1.2pt, colsep=4pt, cells={font=\footnotesize}, row{1}={font=\sffamily\bfseries\footnotesize}}
部品 & 役割 \\
章の帯 & 章番号・題名・部の名前・★。題名は子どもに問いかける形を基本にし、目次にもそのまま載せる \\
道しるべ & 帯のすぐ下。「この章の前に」（2〜3章）と「この章の先に」 \\
できるようになること & その章の目標（普通は3つ）。各項目の頭にその★ \\
はてな？ & 章の問い（身近な不思議）。「数学の道具」の中で必ず答える \\
やってみよう & 実験・観察・電卓・表計算。手順に番号 \\
思い出そう & 前の巻・前の章の要点を2〜4行、章番号付きで。前の巻にあることは繰り返して教えない \\
数学の道具 & 小見出しごとに1つの考え \\
もっと知りたい & 発展。★付き \\
せいかくに言うと & たとえ話のすぐ後に、どこまで正しいかと正確な定義 \\
たいせつ & 1行に1つのこと。定義を置く \\
この章で出てきたことば & 新しい言葉と短い説明 \\
記号カード & 新しい記号・単位の読み方・意味・英語・初めて出た章。1章に2枚まで \\
例題・じぶんでとこう & 「考え方」「式」「答え」の札（第1〜3巻は「しき」「こたえ」）。考え方と途中の式を1行ずつ書く。「じぶんでとこう」は例題と同じ形で同じ数 \\
まちがえやすいところ & ×と○を1行ずつ、もう1行で「なぜ間違いか」 \\
物理・工学ではこう使う & 本物の機械・仕組みの例（第1巻は「つかわれている ところ」）。「大学への種」を添える章もある \\
コラム & 話題の広がり \\
章末問題 & 基本とチャレンジ。問題の形は全て、例題か「じぶんでとこう」で一度見せてある \\
ふりかえり & 「できるようになること」と同じ項目をチェックの形で。各項目に戻るページ \\
ふきだし & ソラ・ミナ・デルタが次の章へつなぐ。1見開きに1回まで \\
専門書をのぞいてみよう & ★★★★以上の章。専門書での書き方と、1行ずつの訳 \\
「ここまでで十分」の旗 & 3段の章の「手でわかる」の終わり（1.5） \\*
\SetCell[c=2]{c} \kgbcapline & \\
\end{bktable}
\par\addvspace{4pt}

\begin{figure}[p]
\noindent\begin{minipage}{\linewidth}\centering
\bkSeriesMap
\bkfigcaption{全9巻の路線図（縦が巻、横が路線。各巻の前付けにあるものと同じ図）}\label{kgd-f-map}
\end{minipage}
\end{figure}

1冊の本は、次の順に並んでいます。

\begin{bkfigure}
\begin{tikzpicture}[x=1mm, y=1mm]
  \foreach \kgdt/\kgdw/\kgdf/\kgdx in {表紙/12/bkVol0!25/0.5, 扉/9/white/13, 目次/11/white/23, 路線図/13/white/35, この本の\\使い方/16/white/49, 序章・部扉・\\章（本文）/31/bkVol0!12/66, 巻末/15/white/98, 奥付/11/white/114, 裏表紙/12.6/bkVol0!25/126} {
    \node[draw=bkVol0, line width=0.8pt, fill=\kgdf, minimum width=\kgdw mm, minimum height=13mm, anchor=west, align=center, font=\sffamily\footnotesize, text=bkInk, inner sep=0.5mm] at (\kgdx,0) {\kgdt};}
  \node[kgb lab, anchor=north west, text=bkGray, text width=138mm, align=left] at (0,-8) {扉は1ページ、目次は2〜3ページ（第1巻は2〜5ページ）、路線図は4〜5ページ（第1巻は6〜7ページ）、「この本の使い方」は6〜7ページ（第1巻は8〜9ページ）。巻末の中身は 1.9。};
\end{tikzpicture}
\bkfigcaption{1冊の並び（表紙の裏には案内役3人の紹介、裏表紙には全巻の一覧）}
\end{bkfigure}

\kgbsub{1.5}{★と3つの段、「名前」と「実際の操作」の型}

各章と「できるようになること」の各項目には、★を1〜5個付けています。★は、このシリーズの中での難しさ、つまり、その章でどの道具を使うかの目安です（★の数と内容の高さのおおよその対応は、9.1 の表）。範囲のある章は「★〜★★」のように書き、「この本の使い方」では★を「どこまで登るか」の山の図で説明しています。子どもには「自分の読める★まで読んで、先へ進んでもだいじょうぶ」と伝えています。

大学の内容を扱う第6〜9巻では、紙面に2つの型を使っています。1つめは章の3段です。大学の内容も中学生の言葉で書くため（編集方針8）、6-10〜6-14、7-7〜7-10、8-1〜8-13、9-1〜9-11 の各章を次の3段に分けています（第6〜9巻の★★★★★の章は、読みものと終章を除いてすべて含まれます）。

\begin{bkfigure}
\begin{tikzpicture}[x=1mm, y=1mm]
  \foreach \kgbi/\kgbic/\kgbn/\kgbd/\kgbs/\kgbf in {0/me/目でわかる/図と実験/3/bkVol0!6,1/te/手でわかる/表計算・プログラム/4/bkVol0!6,2/shiki/式でわかる/専門書の式/5/white} {
    \node[kgb box, minimum width=38mm, minimum height=21mm, text width=36mm, anchor=north west, fill=\kgbf] (s\kgbi) at ({\kgbi*48},0) {};
    \node[anchor=north west, inner sep=0pt] at ({\kgbi*48+2.5},-2.5) {\bkicon[1.6em]{\kgbic}{bkVol0}};
    \node[kgb lab, anchor=north west, font=\sffamily\bfseries\small] at ({\kgbi*48+11},-3.2) {\kgbn};
    \node[kgb lab, anchor=north west] at ({\kgbi*48+2.5},-11) {\kgbd};
    \node[anchor=north west, inner sep=0pt, font=\sffamily\footnotesize] at ({\kgbi*48+2.5},-16) {項目の★：\bkstars{\kgbs}};}
  \draw[kgb arr] (s0) -- (s1);
  \draw[kgb arr] (s1) -- (s2);
  \draw[bkVol0, line width=1pt] (88.5,0.5) -- (88.5,9);
  \fill[bkVol0] (88.5,9) -- (100,6.6) -- (88.5,4.2) -- cycle;
  \node[kgb lab, anchor=west, font=\sffamily\bfseries\footnotesize] at (101,6.6) {ここまでで十分};
  \draw[decorate, decoration={brace, mirror, amplitude=2mm}, bkGray, line width=0.8pt] (0,-23) -- (86,-23);
  \node[kgb lab, anchor=north] at (43,-26.5) {使う道具は★★★まで};
\end{tikzpicture}
\bkfigcaption{章の3段（6-10〜6-14、7-7〜7-10、8-1〜8-13（8-5 は2段）、9-1〜9-11）}\label{kgd-f-dan}
\end{bkfigure}

「目でわかる」（図と実験）と「手でわかる」（表計算・プログラム）は★★★以下の道具だけでたどれ、「手でわかる」の終わりに「ここまでで十分」の旗を立てています。その後「式でわかる」で★★★★★の式を作ります。「できるようになること」の3つの項目は、普通、段ごとに★★★・★★★★・★★★★★と並びます。図\ref{kgd-f-dan}の「項目の★」は、その段で到達する内容の高さ、「使う道具は★★★まで」は、そこまでたどるのに使う計算の高さを表します。

2つめは、「名前」と「実際の操作」の型です。$\iint$ や $\mathcal{L}[f(t)]$ のような記号は「操作の名前」で、記号を見ただけではまだ計算できません。実際の値は、第8巻では「ふつうの定積分（第5巻）の繰り返し」に、第9巻では4つの道具（ふつうの定積分・表を引く・部分分数・行列の計算）のどれかに直して求めます。紙面では「名前の段→直す段→楽をする段」の3つの段、場合分けの表、流れ図を使い、札の形（点線の枠・実線の枠・塗りつぶし）で段を見分けられるようにしています。これは、依頼者からの「操作の名前と実際の操作の違いは相当わかりにくく、どのような操作をすれば求められるのかも場合によりけりなので、この場合はどうすれば良いのかを丁寧にする必要がある」という指示（10月29日）を紙面にしたもので、章の3段とは別の、その1つ下の段に当たります。

\par\addvspace{4pt}
\kgbcap{第6巻からの2つの型（「名前」と「実際の操作」の型は、第6・7巻にはありません）}{kgd-t-kata}
\begin{bktable}[long, theme=kgb]{hline{Z}={0pt,white}, vline{1-Z}={Z}{0pt,white}, row{Z}={abovesep=3pt, belowsep=1pt}, hspan=minimal, rowhead=1, colspec={Q[l,m,wd=34mm] Q[l,m,wd=56mm] Q[l,m,wd=36mm]}, rowsep=2pt, colsep=4pt, cells={font=\small}, row{1}={font=\sffamily\bfseries\small}}
型 & 中身 & 使っている巻・章 \\
章の3段 & 「目でわかる」→「手でわかる」→「式でわかる」。「手でわかる」の終わりに「ここまでで十分」の旗 & 第6〜9巻：6-10〜6-14、7-7〜7-10、8-1〜8-13（8-5 は2段）、9-1〜9-11 \\
「名前」と「実際の操作」 & 名前の段→直す段→楽をする段の3つの段、場合分けの表、流れ図 & 第8巻（8-0 予告・8-6〜8-11・8-14・巻末の「積分の地図」）、第9巻（全章と巻末の「変換の地図」） \\*
\SetCell[c=3]{c} \kgbcapline & & \\
\end{bktable}
\par\addvspace{4pt}

\kgbsub{1.6}{案内役：ソラ・ミナ・デルタ}

どの巻にも、3人の案内役が登場します。表紙の裏での紹介の他、序章・部扉と、各章のふきだし（1見開きに1回まで）に出てきて、問いかけたり、次の章へつないだりします。

\begin{bkfigure}
\begin{tikzpicture}[x=1mm, y=1mm]
  \bkChar{sora}{(18,0)}{30mm}\bkChar{mina}{(62,0)}{30mm}\bkChar{delta}{(108,0)}{19mm}
  \node[kgb lab, anchor=north, font=\sffamily\bfseries\small] at (18,-1.5) {ソラ};
  \node[kgb lab, anchor=north, font=\sffamily\bfseries\small] at (62,-1.5) {ミナ};
  \node[kgb lab, anchor=north, font=\sffamily\bfseries\small] at (108,-1.5) {デルタ};
  \node[kgb lab, anchor=north, text width=40mm, align=center] at (40,-7.5) {子ども。巻が進むごとに成長し、\\第7巻からは大学生};
  \node[kgb lab, anchor=north, text width=40mm, align=center] at (108,-7.5) {三角形のロボット。\\名前は $\Delta$（デルタ）から};
\end{tikzpicture}
\bkfigcaption{3人の案内役}
\end{bkfigure}

\kgbsub{1.7}{どこから始めるか}

学年の目安（表\ref{kgd-t-vols}）は、読む時期を決めるものではありません。始める所は、次の決まりで選びます。

\begin{taisetsu}
いま学校で習っている学年の内容が含まれる、一番前の巻・章から始める
\end{taisetsu}

学年の内容がどの巻・章にあるかは、表\ref{kgb-t-gakunen}（4.3）によります。高校の数学は、1つの科目の内容がいくつもの巻に分かれているので、いま習っている単元が含まれる一番前の章から始めます。始めた章より前の章にも、まだ習っていない内容が少し含まれることがあります。分からない所は、章の「道しるべ」の「この章の前に」の章へ戻ってください。巻の順に読むのが基本ですが、自分の読める★まで読んで先へ進んでも構いません。子ども用の『はじめに読む本』の流れ図「はじめる巻を決める」も、同じ決まりで作っています。

\par\addvspace{4pt}
\kgbcap{いま習っていることと、始める巻・章}{kgd-t-start}
\begin{bktable}[long, theme=kgb]{hline{Z}={0pt,white}, vline{1-Z}={Z}{0pt,white}, row{Z}={abovesep=3pt, belowsep=1pt}, hspan=minimal, rowhead=1, colspec={Q[l,m,wd=27mm] Q[l,m,wd=33mm] Q[l,m,wd=66mm]}, rowsep=1.6pt, colsep=4pt, cells={font=\footnotesize}, row{1}={font=\sffamily\bfseries\footnotesize}}
いま習っていること & 始める巻・章 & その学年の内容がある主な章と、進め方 \\
小学1年 & 第1巻の序章（1-0） & 小1の内容は第1巻だけ。第1巻は本に直接書き込んで使います。大人が一緒に読み、「やってみよう」の実験を手伝ってください \\
小学2年 & 第1巻 1-2 & 小2の内容（4けたの数、筆算、九九、長さとかさ）は 1-2 からの第1巻 \\
小学3年 & 第1巻 1-2 & 小3の内容（万の位、3けたの筆算、割り算、小数と分数の入口）は 1-2 からの第1巻。掛け算の筆算は学校の教科書で（7.2） \\
小学4年 & 第1巻 1-15 & 小4の内容は 1-15（余りのある割り算の確かめの式）・1-18（折れ線グラフ）・1-20（角）・1-21（面積）と、第2巻。割り算の筆算は学校の教科書で（7.2） \\
小学5年 & 第2巻の序章（2-0） & 小5の内容は第2巻（2-1 から）と第4巻（図形）。第2巻からは問題をノートに解きます \\
小学6年 & 第2巻 2-4 & 小6の内容（分数の掛け算・割り算、文字、比、比例と反比例）は 2-4 から。2-1〜2-3 にも、まだ習っていない内容（負の数 2-2 など）があるので、道しるべで戻って読みます \\
中学1年 & 第2巻 2-2 & 中1の内容は第2巻（正負の数 2-2 から、文字式・一次方程式・比例と反比例）、第3巻（座標・関数・データ）、第4巻（平面図形・空間図形） \\
中学2年 & 第2巻 2-12 & 中2の内容は第2巻（式の計算 2-12・等式の変形 2-14）、第3巻（連立方程式・一次関数・箱ひげ図）、第4巻（平行線と角・合同） \\
中学3年 & 第2巻 2-6 & 中3の内容は第2巻（誤差 2-6・平方根 2-8・無理数 2-9）、第3巻（展開・二次方程式・$y=ax^2$）、第4巻（相似・三平方の定理） \\
高校の数学Ⅰ & 第2巻 2-9 & 数学Ⅰの内容は第2巻（実数 2-9）、第3巻（二次関数・データの分析）、第4巻（三角比）。第2巻の表計算（2-18）・プログラム（2-19）は、後の巻で使います \\
高校の数学Ⅱ & 習っている単元の最初の章（指数・対数は 2-7、三角関数は 4-2、微分・積分は第5巻の序章 5-0） & 数学Ⅱの内容は第2〜6巻と第9巻に分かれています（表\ref{kgb-t-gakunen}）。数学Ⅱの微分・積分を習っていて、数学Ⅲはまだの人は、第5巻から（数学Ⅲの内容も含みます） \\
高校の数学Ⅲ & 第5巻 5-3 & 極限・微分法・積分法は第5巻（5-3 から）。合成関数・逆関数は第3巻（3-3・3-12） \\
大学生・社会人の学び直し & 学び直したい科目の巻（表\ref{kgb-t-univ}） & 微分・積分は第5巻、微分方程式は第6巻、行列は第7巻。前提になる章は、各章の道しるべでたどれます \\
工学の専門書を読むため & 第6巻（全巻の乗り換え駅） & 第6〜9巻の3段の章は、「ここまでで十分」の旗まで読んでから「式でわかる」へ。第8巻・第9巻の「名前」と「実際の操作」の型で、専門書の記号を計算に直せるようになります \\*
\SetCell[c=3]{c} \kgbcapline & & \\
\end{bktable}
\par\addvspace{4pt}

\kgbsub{1.8}{電卓・表計算・プログラム（Python）}

このシリーズは「計算の量より意味」を大切にしています（編集方針5）。手で計算する練習は第1巻（九九と、足し算・引き算の筆算）で行い、第2巻からは電卓・表計算・プログラムに任せてよいことにしています。掛け算・割り算の筆算は今の版にないので、学校の教科書で練習してください（7.2）。各巻の「この本の使い方」にある「3つのやくそく」は、次の3つです。

\kgbitem{数にはいつも単位}
\kgbitem{式を書く前に表や図}
\kgbitem{くり返しの計算は電卓やコンピュータにまかせてよい}

\par\addvspace{4pt}
\kgbcap{電卓・表計算・プログラムを使い始める章}{kgd-t-tools}
\begin{bktable}[long, theme=kgb]{hline{Z}={0pt,white}, vline{1-Z}={Z}{0pt,white}, row{Z}={abovesep=3pt, belowsep=1pt}, hspan=minimal, rowhead=1, colspec={Q[l,m,wd=24mm] Q[l,m,wd=20mm] Q[l,m,wd=82mm]}, rowsep=2pt, colsep=4pt, cells={font=\footnotesize}, row{1}={font=\sffamily\bfseries\footnotesize}}
道具 & 使い始め & 内容 \\
電卓 & 2-1 から & 四則と $\sqrt{\phantom{2}}$ のキー（2-8）までは普通の電卓でできます。整数でない指数の \texttt{x\^{}y} キー（3-11）、\texttt{log} キー（3-12）、三角比と角の単位の切りかえ（4-2・4-6）からは関数電卓（または表計算）を使います \\
表計算 & 2-18 で学ぶ & セルと番地、「＝」で始まる式、式のコピー、\$ で固定する番地、上の行の結果で下の行を計算する繰り返しの計算、表からグラフ。以後の全巻で使い、第6巻のオイラー法（6-2）もこの形で解きます \\
プログラム & 2-19 で学ぶ & ブロックをつないで書く方法（Scratch など）と文字で書く方法（Python など）を並べて紹介。変数と代入、繰り返し、条件分岐、小数のずれ \\
Python & 第7巻・第8巻の本文など & 行列の計算（7-1：表計算の MMULT と Python）、2変数の関数の地図（8-1：表計算と Python）など。第6〜9巻の巻末には「プログラムのひな形」 \\*
\SetCell[c=3]{c} \kgbcapline & & \\
\end{bktable}
\par\addvspace{4pt}

第1巻は本に直接書き込んで使います。第2巻からは、問題をノートに解きます。ノートを1冊用意してください。

\kgbsub{1.9}{答えと巻末}

巻末の「こたえ」は、「じぶんでとこう」と章末問題の答えだけです（考え方は、今後作る問題集と指導書に載せる予定です）。最後に「答え合わせのしかた」があります。巻末には、他に次のものがあります。

\par\addvspace{4pt}
\kgbcap{巻末の中身}{kgd-t-back}
\begin{bktable}[long, theme=kgb]{hline{Z}={0pt,white}, vline{1-Z}={Z}{0pt,white}, row{Z}={abovesep=3pt, belowsep=1pt}, hspan=minimal, rowhead=1, colspec={Q[l,m,wd=38mm] Q[l,m,wd=88mm]}, rowsep=1.5pt, colsep=4pt, cells={font=\footnotesize}, row{1}={font=\sffamily\bfseries\footnotesize}}
巻末の部品 & 中身 \\
こたえ & 「じぶんでとこう」と章末問題の答え。最後に「答え合わせのしかた」 \\
記号カード & 各章の記号カードを集めた一覧（第4巻からはギリシャ文字の読み方、記号を書くときの注意なども） \\
よく使う定数・単位・公式 & 第1巻は「たんいのひょう」。第6巻からは「この巻でくり返し使う場面と数」と「プログラムのひな形」、第8巻は初めに見開きの「積分の地図」、第9巻は初めに「変換の地図」 \\
チェック表 & 各章の「ふりかえり」を集めた表 \\
つぎの巻へ & 次の巻の予告と「この巻の道具は、この先でこう使う」。第9巻は代わりに「物理シリーズへ」 \\
さくいん & その巻の索引。第9巻には第1〜9巻を通した「シリーズ全体のさくいん」 \\
ふろく（第1巻だけ） & コピーして使う10のわく・かずカード・ものさし・とけい・方眼 \\*
\SetCell[c=2]{c} \kgbcapline & \\
\end{bktable}
\par\addvspace{4pt}

\kgbsub{1.10}{問題集と指導書（今後の予定）}

このシリーズは教科書を優先して作りました。問題集と、教員向けの指導書は、今後作る予定で、まだ刊行していません。巻末の「こたえ」に載せていない問題の考え方や、章ごとの学年の対応は、指導書に載せる予定です。それまでの間、学年との対応は、この本の第4章と付録A・付録Bをご覧ください。

\kgbsub{1.11}{物理シリーズへ}

このシリーズの後には、物理シリーズ（全5巻）が続きます。書き方は数学シリーズの第4巻以降と同じ（ルビなし・常用漢字まで）で、各章の道しるべに、前に学ぶ数学の章の番号（6-2 など）を書きます。数学シリーズの第9巻の巻末「物理シリーズへ」では、物理の「なぜ？」とこのシリーズの章の対応表と、物理5巻の紹介を載せています。

\par\addvspace{4pt}
\kgbcap{物理シリーズ（全5巻）と、前に学ぶ数学の巻}{kgd-t-phys}
\begin{bktable}[long, theme=kgb]{hline{Z}={0pt,white}, vline{1-Z}={Z}{0pt,white}, row{Z}={abovesep=3pt, belowsep=1pt}, hspan=minimal, rowhead=1, colspec={Q[c,m,wd=16mm] Q[l,m,wd=30mm] Q[l,m,wd=78mm]}, rowsep=1.5pt, colsep=4pt, cells={font=\small}, row{1}={font=\sffamily\bfseries\small}}
物理の巻 & 巻名 & 前に学ぶ数学の巻（章ごとの指定から集計） \\
1 & 力学 & 第2〜8巻 \\
2 & 熱学 & 第2・3・5〜9巻 \\
3 & 波動 & 第3〜9巻 \\
4 & 電磁気学 & 第3〜9巻 \\
5 & 原子 & 第2〜9巻 \\*
\SetCell[c=3]{c} \kgbcapline & & \\
\end{bktable}
\par\addvspace{4pt}
